% ===================== METODOLOGÍA =====================
\section{Metodología}

\subsection{Propiedades Termodinámicas del Agua/Vapor}
Las entalpías específicas del agua y vapor se obtuvieron de las formulaciones IAPWS-97 (International Association for the Properties of Water and Steam), implementadas computacionalmente mediante la librería \texttt{iapws} de Python.

\subsection{Balance de Materia}
El balance de materia en el lado agua de la caldera considera la producción neta de vapor, la purga continua del domo y el flujo total de agua de alimentación:
\begin{equation}
	m_{fw} = m_{stm} + m_{purge} \quad \text{donde} \quad m_{purge} = 0.02 \cdot m_{fw}
\end{equation}

\subsection{Balance de Energía}
El calor absorbido por el fluido de trabajo se calcula como:
\begin{equation}
	Q_{abs} = m_{stm} \cdot h_{stm} + m_{purge} \cdot h_{purge} - m_{fw} \cdot h_{fw}
\end{equation}

El calor requerido del combustible, considerando la eficiencia de la caldera ($\eta = 94\%$):
\begin{equation}
	Q_{fuel} = \frac{Q_{abs}}{\eta}
\end{equation}

El consumo de bagazo se determina dividiendo la energía requerida entre el PCI:
\begin{equation}
	m_{bagazo} = \frac{Q_{fuel}}{PCI}
\end{equation}

\subsection{Poder Calorífico}
Se empleó la fórmula de Boie (Ecuación~\ref{eq:boie}) para el cálculo del PCS a partir de la composición elemental, por ser una de las correlaciones más precisas para biomasas según la literatura especializada \cite{boie1953}.
